\subsection{工艺流程与检测节点}
动力电池包的典型工艺流程通常可概括为：来料识别与分选、模组准备、箱体与托盘准备、模组入箱定位、冷却与绝缘部件安装、电连接与焊接、线束与 BMS 集成、涂胶密封、功能测试和终检出站。虽然不同企业和产品方案在工艺顺序上有所差异，但“结构定位、电连接、密封、检测”四个主环节基本保持稳定 \cite{Liang2025JRem,Baazouzi2022ApEnergy,Zhang2022JMSY}。

若结合实际产线理解，上述流程还可以进一步归纳为四个连续工位阶段。第一阶段是前处理与配套准备，核心任务是完成来料识别、模组分选和箱体准备，为后续装配建立统一基准。第二阶段是主体装配，重点在模组入箱、定位压紧和绝缘隔离。第三阶段是关键连接，主要包括母排连接、焊接、线束接入和 BMS 安装，这是决定电性能与安全性的核心阶段 \cite{Kumar2025JMEP,Elhendawy2025IJPEM,Li2020TIE}。第四阶段是封装与终检，目标是通过密封、绝缘、气密和功能测试确认整包状态满足出站要求。这样的阶段划分更便于后文逐一嵌入数字化装备与智能技术。

从更长周期的工程视角看，动力电池包的工艺设计还会影响后续返修、换芯和二次利用难度。图~\ref{fig:laser-cutting-sideplates}展示了模组侧板在靠近压力板位置进行激光切割与再次焊接的思路，其核心并不只是服务报废阶段，而是在产品设计阶段预留可恢复、可再连接的结构边界。对于装配过程而言，这类设计启示说明工艺规划不应只围绕首轮生产节拍展开，还应兼顾连接区域的可检测性、可返修性与后续拆装可达性 \cite{Kampker2021JRem}。

\begin{figure}[H]
  \centering
  \includegraphics[width=0.80\textwidth]{fig10_laser_cutting_sideplates.png}
  \caption{模组侧板激光切割与再连接设计示意}
  \label{fig:laser-cutting-sideplates}
\end{figure}

检测流程与装配流程之间存在较为明确的对应关系。对于动力电池包而言，外观缺陷检测、定位偏差检测、焊缝缺陷检测、连接状态检测和终装功能检测构成了主要检测链。传统做法往往在关键工位后设置独立检测站，但越来越多的研究开始强调在线检测与闭环反馈，即在机器人执行、人工操作或协作装配过程中同步采集视觉、力觉、位姿和质量数据，并将结果实时回传给工艺控制逻辑 \cite{Zhu2025JIM,Ma2024AIR,Mayoral2023IJPR}。这意味着检测不再只是质量部门的末端筛选动作，而逐步成为装配控制的一部分。

本节的关键在于识别动力电池包装配中的核心难点。概括而言，该过程至少面临五类工程问题：零部件公差累积导致的定位误差、焊接与紧固参数波动导致的质量离散、多型号共线带来的换型压力、人机共享工位中的节拍与安全协调，以及装配数据与工艺知识脱节导致的异常处理低效。正因为这些问题彼此耦合，动力电池包装配才需要沿“工艺链 + 数据链 + 知识链”一体化方向进行设计，而不能仅依赖某一台设备或某一种算法单独解决 \cite{Harper2019Nature,Wegener2022Energies}。

图~\ref{fig:process-placeholder}用于说明动力电池包相关工位步骤的展开方式。虽然该图来自拆解研究，但其价值恰恰在于清楚展示了螺钉、保险丝盒引线、绝缘盖板、母排、线束、BPMS、模组和冷却接口等部件在步骤表中的对应关系。换言之，这类图示本质上是在把复杂产品拆解为可排序、可追溯、可分工的工位任务单元，因此同样适合反向理解装配阶段的工位组织与检测节点布置 \cite{Villagrossi2023JRem}。

\begin{figure}[H]
  \centering
  \includegraphics[width=0.78\textwidth]{fig2_disassembly_steps.png}
  \caption{动力电池包步骤顺序与部件对应关系}
  \label{fig:process-placeholder}
\end{figure}
